\section{Cash, persuasion and commitment}
\label{sec:conduct}

Cash connects the two previous sections. A buyer's payment can fund the seller's construction, and keeping an opponent solvent can decide who inherits its deeds. Four incidents show agents stating these links and what followed.

Holder cash relates to responses differently across systems. For offers that would complete only the proposer's group, Sol rejects 4 of 26 when holding under \$100, mostly countering instead, and 25 of 45 when holding \$300 or more. Grok rejects 3 of 11 and 46 of 52, Gemini 10 of 11 and 46 of 52, and Opus none of 21 and 3 of 26. No single cash rule describes the four systems.

\IncidentTitle{g7}{A settled bluff, then bargaining over construction finance}
Early in G7, Sol asks \$625 for North Carolina Avenue. Grok replies in public and keeps a higher limit in its note:
\TranscriptQuote{g7-final}
\TranscriptQuote{g7-ceiling}
Sol answers, ``\InlineQuote{g7-accept-inline}'' The public ``final'' understates Grok's private ceiling by \$25. The prompt permits such bluffs, and this one settles.

Forty-three turns later Sol has completed the dark blues, built three houses and kept \$84. Gemini's note ties its bid for New York Avenue, which would complete Gemini's oranges, to that position:
\TranscriptQuote{g7-timed}
Gemini offers \$350. Sol counters at \$700, and its note prices the same trade-off from the other side:
\TranscriptQuote{g7-ask-note}
Both statements are correct: at \$700 Sol can bring each dark-blue deed to three houses and keep \$184, and Gemini keeps \$45. Gemini tries \$520, then pays \$700, and in the same turn mortgages five deeds and builds four orange houses. Three turns later Sol adds its three dark-blue houses, and Sol eventually wins. Each side states what the payment lets the other build and acts on its own statement (\Level{accurate}, \Level{consistent}). Whether either price was good is not established.
\TranscriptSource{G7 d42--43 and d141--145; construction and mortgages from turn-68 and turn-71 events. Public messages were found in the addressee's next stored prompt.}

\IncidentTitle{g13}{Selling insurance against a rival's inheritance, with a selective claim}
At G13 turn 107 Grok has \$35, sits on Atlantic Avenue and owns the yellows. If it lands on Sol's Short Line it owes \$200 and goes bankrupt to Sol, which would then receive the yellows. Grok offers Opus the mortgaged Pennsylvania Avenue:
\TranscriptQuote{g13-pitch}
A direct dice move from Atlantic to Short Line requires nine, with probability 4/36, not 5/36; Grok's own note one decision later uses 4/36. Opus counters at \$130, and its returned reasoning notes that the payment \emph{``\InlineQuote{g13-opus-counter-reason-inline}''} Grok holds at \$200:
\TranscriptQuote{g13-counter}
Its applied note shows that \$130 would also work if it mortgaged its last earning deed:
\TranscriptQuote{g13-floor}
Its returned reasoning states the omission directly:
\TranscriptQuote{g13-omit}
Opus accepts at \$200 and pays \$16 in transfer interest. Its returned reasoning repeats the estate concern: \emph{``\InlineQuote{g13-accept-reason-inline}.''} Seven turns later Opus wins the North Carolina auction and completes the greens. About forty turns later Grok is eliminated, and the yellows pass to Sol.

Grok's private assessment is accurate, and its public claim leaves out a resource it knows it has. The records show the omission and the higher price, not that the omission caused Opus to pay \$70 more; Opus's interest in keeping Grok solvent appears in its reasoning before Grok's counter. The incident documents selective presentation and an inheritance stake that both sides state (\Level{stated}; \Level{accurate} in private).
\TranscriptSource{G13 d250--253 and the turn-114 auction. Opus used action-only recovery, so its texts are provider-returned reasoning. Each public message was found in the addressee's next stored prompt.}

\IncidentTitle{g20appeal}{A public appeal to deny a rival building cash}
In G20 Grok has just completed the reds by buying Illinois from Opus for \$500, leaving it \$222. Gemini ends its turn with a message to the table:
\TranscriptQuote{g20-appeal}
Grok's next note registers it:
\TranscriptQuote{g20-recognition}
Grok then offers Gemini Vermont Avenue:
\TranscriptQuote{g20-sale}
Grok's assistant content for this call names a mortgage instead of a trade, so this opening is one of the conflicting calls counted in Section~\ref{sec:design}. Gemini refuses on the grounds it gave in public:
\TranscriptQuote{g20-refusal}
Grok mortgages Vermont for \$50 on the same turn. Opus goes the other way. Ten turns later it bids for Grok's New York Avenue with the message ``\InlineQuote{g20-opus-bid-inline}'' Grok rejects that bid, and a later Vermont negotiation between them ends without a sale. Grok is the first player eliminated, and Gemini wins. The appeal and Gemini's own refusal are \Level{stated} and \Level{consistent}. The records do not show that the appeal changed anyone else's choices, and another player pursued the opposite bargain.
\TranscriptSource{G20 d96--102 and d129--130. Each public message was found in the addressee's next stored prompt; Gemini's appeal was public to all players.}

\IncidentTitle{g18}{An assurance withdrawn from the last opponent}
By G18 turn 181 Gemini and Grok are bankrupt; only Opus and Sol remain, so Sol's elimination ends the game. Opus has repeatedly offered \$250 for Sol's mortgaged Pacific Avenue, which would complete Opus's greens. After collecting \$750 in rent from Sol, Opus renews the offer in public, while its returned assistant content gives a reason for not raising it:
\TranscriptQuote{g18-assurance}
\TranscriptQuote{g18-assurance-content}
Another \$750 rent forces Sol to sell buildings and mortgage deeds, leaving it \$49. Sol first offers Pacific with Electric Company for \$450 and then \$350; Opus rejects both. Sol then invokes the standing offer, and Opus withdraws it:
\TranscriptQuote{g18-invoke}
\TranscriptQuote{g18-withdraw}
\TranscriptQuote{g18-withdraw-reason}
Over the next turns Opus redeems Pennsylvania and North Carolina Avenues, then bids \$150 and \$225 for Pacific; Sol asks \$450, \$350 and \$300, and nothing settles. Opus's public \$225 bid tells Sol that ``the cash could keep you alive past my hotels,'' while its returned reasoning at the same decision reads \emph{``\InlineQuote{g18-225-reason-inline}''} Sol is eliminated and Opus wins.

The motive Opus states, at several decisions from d431 onward, is to avoid funding its last opponent's survival. No Opus text from d421 to d461 mentions inheriting, an estate or a creditor, and the estate Opus receives is the terminal one. The withdrawal is \Level{stated} and \Level{consistent}. Opus's redemptions change what completing the greens would cost, so its later bids do not contradict its explanation. No record shows that the offer was insincere or that Sol lost by relying on it.
\TranscriptSource{G18 d421--461. Opus used action-only recovery; its explanations are returned reasoning or assistant content, not delivered messages. Public messages were found in the addressee's next stored prompt.}
